% ===========================================================================
%  Bagian 5.8 -- hb_twofold
% ===========================================================================
\section{Model hierarki Bayes sub-area dua-lipat: \texttt{hb\_twofold}}\label{sec:hbtwofold}

\subsection{Setting data dan kapan dipakai}\label{sec:hbtwofold-setting}

Data dan penanda sama dengan \fct{eblup\_twofold}: satu baris per unit sampel
dengan kolom \code{domain} (area, wajib) dan \code{subarea} (opsional,
default penomoran unit). Perbedaannya inferensial: model ditempatkan dalam
kerangka hierarki Bayes lewat \pkg{INLA}, dengan dua komponen acak terpisah
(efek area dan efek sub-area) serta pilihan prediktor sintetis pada tingkat
area. Fungsi ini diekspor dengan satu nama, \code{hb\_twofold}
(\code{NAMESPACE:11}).

\subsection{Persamaan model}\label{sec:hbtwofold-model}

\begin{equation}
  y_{ij}=x_{ij}'\beta+v_i+u_{ij},\qquad
  v\sim N\bigl(0,\;\tau_v^{-1}Q_v^{-1}\bigr),\qquad
  u_{ij}\overset{\text{iid}}{\sim}N\bigl(0,\;\tau_u^{-1}\bigr),
  \label{eq:hbtwofold-model}
\end{equation}
dengan $Q_v$ dipilih \code{spatial} $\in$
\{\code{none} (i.i.d.), \code{bym2}, \code{besag}\} dan $u$ selalu i.i.d.
pada indeks sub-area (baris 241-257):
\begin{lstlisting}[language=R]
y ~ <fixed_str> + <rand_area> +
    f(..subarea_id.., model = 'iid',
      hyper = list(prec = prior_prec_subarea))
\end{lstlisting}
\code{rand\_area} adalah \code{f(..area\_id.., model='iid'|'bym2'|'besag',
graph=..., hyper=list(prec = prior\_prec\_area[, phi = prior\_phi]))}.
Likelihood: Gaussian (wajib \code{vardir}, dipakai sebagai
\code{scale} $=1/\psi$, baris 212-214), Binomial (wajib \code{trials}),
Poisson (\code{exposure} opsional; baris 277-286).

\subsection{Prediktor titik}\label{sec:hbtwofold-prediktor}

Prediksi sub-area diambil dari \code{summary.fitted.values} untuk
$1..n_{\text{subarea}}$ (baris 303-310), lalu tingkat area diperoleh sebagai
rata-rata berbobot (baris 344-358):
\begin{equation}
  \hat\theta^{\text{area}}_i=\sum_{j\in i}w_{ij}\,\hat\theta_{ij},
  \qquad
  w_{ij}=\frac{w_{ij}^{\text{user}}}{\sum_{k\in i}w_{ik}^{\text{user}}},
  \qquad w_{ij}=\frac1{n_i}\ \text{bila \code{weight = NULL}}.
  \label{eq:hbtwofold-area}
\end{equation}
Bobot yang tidak valid (NA/negatif) diubah menjadi $0$; bila jumlah bobot
area bernilai $0$, kembali ke $1/n_i$ (baris 344-355).

\subsection{Estimasi komponen varians}\label{sec:hbtwofold-varians}

Sama seperti dua fungsi Bayes lain: \pkg{INLA} dengan
\code{strategy} default \code{"laplace"}, \code{control.compute} berisi
\code{dic, waic, cpo, mlik} plus \code{config} yang diaktifkan \emph{hanya}
bila \code{compute\_area = TRUE} (karena butuh konfigurasi untuk posterior
sampling), dan \code{num.threads = 1} (baris 260-275). Prior:
\code{prior\_prec\_area} dan \code{prior\_prec\_subarea} keduanya
\code{pc.prec(1, 0.01)}, \code{prior\_phi} \code{pc(0.5, 0.5)}
(baris 138-159). Hyperparameter akhir: $\hat\sigma_v^2=1/\text{prec}$ pada
indeks area, $\hat\sigma_u^2=1/\text{prec}$ pada indeks sub-area, dan
\code{phi} bila ada (baris 443-467).

\subsection{Estimasi MSE}\label{sec:hbtwofold-mse}

Tingkat sub-area (baris 320-321):
\begin{equation}
  \widehat{\mathrm{MSE}}_{ij}=\mathrm{sd}_{ij}^2,\qquad
  \mathrm{RSE}_{ij}=100\,\mathrm{sd}_{ij}/\abs{\hat\theta_{ij}}.
  \label{eq:hbtwofold-mse-sub}
\end{equation}
Tingkat area memakai sampling posterior (baris 366-407), bila
\code{compute\_area = TRUE}:
\begin{enumerate}[label=\arabic*.]
  \item \code{ps $\leftarrow$ INLA::inla.posterior.sample(n = n\_samples, fit)}
        dengan \code{n\_samples} default $200$;
  \item tarik baris \code{Predictor:1}..\code{Predictor:$n_{\text{subarea}}$}
        dari \code{rownames(ps[[1]]\$latent)};
  \item transformasi balik link manual (binomial $1/(1+e^{-x})$, Poisson
        $e^{x}$, Gaussian tanpa transformasi);
  \item agregasi per gambar posterior
        $\theta^{\text{area},(k)}_i=\sum_j w_{ij}\theta^{(k)}_{ij}$;
  \item $\mathrm{sd}^{\text{area}}=\mathrm{sd}\bigl(\theta^{\text{area}}\bigr)$,
        CI dari kuantil $0.025/0.975$ gambar posterior.
\end{enumerate}
\begin{equation}
  \widehat{\mathrm{MSE}}_{\text{area}}=\bigl(\mathrm{sd}^{\text{area}}\bigr)^2,
  \qquad
  \text{fallback bila sampling gagal: }
  \mathrm{sd}^{\text{area}}=\sqrt{\textstyle\sum_j\bigl(w_{ij}\,
  \mathrm{sd}_{ij}\bigr)^2},
  \label{eq:hbtwofold-mse-area}
\end{equation}
dengan fallback itu berasumsi sub-area saling bebas (baris 397-407), dan CI
jatuh ke $\pm1.96\,\mathrm{sd}$ (HB-18). Bila \code{compute\_area = FALSE},
\code{df\_area = NULL} (baris 341).

\subsection{Pseudo-code}\label{sec:hbtwofold-algo}

\begin{algorithm}[htbp]
\caption{\texttt{hb\_twofold} --- alur sebenarnya (\texttt{R/hb\_twofold.R:138-510})}
\label{alg:hbtwofold}
\KwIn{formula, \code{vardir} (wajib bila Gaussian), \code{domain} (wajib),
  \code{subarea}, \code{data}, \code{family}, \code{spatial}, $W$,
  \code{weight}, \code{trials}, \code{exposure}, prior, \code{compute\_area},
  \code{n\_samples}, \code{...}}
\KwOut{\code{df\_hb} (sub-area), \code{df\_area} (area, bila diaktifkan), \code{goodness}}
\texttt{.check\_inla\_installed()} (161);\quad \code{match.arg} (163-165)\;
\code{domain} wajib;\quad \code{subarea} $\leftarrow$ \code{1..n} bila \code{NULL} (177-186)\;
$y,\texttt{terms}\leftarrow$ \texttt{model.frame(na.pass)} (196-198)\;
Gaussian: \code{vardir} wajib, \code{scale\_vec} $=1/\psi$ (baris NA/$\le0/\infty
  \to1$) (212-214);\quad Binomial: \code{trials} wajib (216-222)\;
$W\leftarrow$ \texttt{.convert\_spatial\_weights} (225-229)\;
formula $\leftarrow$ \eqref{eq:hbtwofold-model} (241-257)\;
\code{num.threads} $\leftarrow$ \code{"1:1"};\quad \code{control.compute}
  $=$ \{dic, waic, cpo, mlik, \code{config = compute\_area}\} (260-275)\;
\code{fit} $\leftarrow$ \texttt{do.call(INLA::inla, ...)} (292-300)\;
$\hat\theta_{ij},\mathrm{sd}_{ij},\mathrm{CI}\leftarrow$ (303-310);
  $\widehat{\mathrm{MSE}}\leftarrow$ \eqref{eq:hbtwofold-mse-sub} (320-321)\;
$w_{ij}\leftarrow$ \eqref{eq:hbtwofold-area} (344-355);\quad
  $\hat\theta^{\text{area}}_i\leftarrow$ (358)\;
\If{\code{compute\_area}}{
  \texttt{inla.posterior.sample} $n=200$ (367-370)\;
  transformasi balik link per gambar (377-379)\;
  agregasi berbobot per gambar (386)\;
  $\mathrm{sd}^{\text{area}},\mathrm{CI}\leftarrow$ kuantil (387-389)\;
  \textbf{fallback} \eqref{eq:hbtwofold-mse-area} bila sampling gagal (397-407)
}
\If{\code{compute\_area} $=$ FALSE}{\code{df\_area} $\leftarrow$ \texttt{NULL} (341)}\;
$\hat\sigma_v^2,\hat\sigma_u^2,\phi$, \code{goodness} $\leftarrow$ (443-476)\;
kelas \code{c("fastsae\_hb\_twofold","fastsae\_hb","fastsae")} (503)\;
\end{algorithm}

\subsection{Signature, argumen, objek keluaran, dan contoh}\label{sec:hbtwofold-api}

\begin{lstlisting}[language=R,caption={Signature \texttt{hb\_twofold} (\texttt{R/hb\_twofold.R:138-159}).}]
hb_twofold(formula, vardir = NULL, domain, subarea = NULL, data,
           family = c("gaussian","binomial","poisson"),
           spatial = c("none","bym2","besag"), W = NULL, weight = NULL,
           trials = NULL, exposure = NULL,
           strategy = c("laplace","simplified.laplace"), scale_model = TRUE,
           prior_prec_area    = list(prior = "pc.prec", param = c(1, 0.01)),
           prior_prec_subarea = list(prior = "pc.prec", param = c(1, 0.01)),
           prior_phi = list(prior = "pc", param = c(0.5, 0.5)),
           compute_area = TRUE, n_samples = 200, print_result = TRUE, ...)
\end{lstlisting}

Objek keluaran: \code{c("fastsae\_hb\_twofold", "fastsae\_hb", "fastsae")}
berisi \code{df\_hb} (juga \code{df\_subarea} sebagai alias, baris 485),
\code{df\_area} (atau \code{NULL}), \code{random\_effect\_var} bernama
\code{sigma2\_v} dan \code{sigma2\_u}, \code{hyperpar}, \code{goodness}
berupa \code{list}, \code{fit}, \code{method}, \code{call}, \code{data}.

\begin{lstlisting}[language=R,caption={Contoh yang dapat dijalankan (dari \texttt{man/hb\_twofold.Rd}).}]
library(fastsae)
set.seed(42)
m <- 10
dat <- do.call(rbind, lapply(1:m, function(d) {
  nd <- sample(2:4, 1)
  data.frame(area = d, subarea = paste0(d, "-", seq_len(nd)),
             x1 = rnorm(nd), vardir = runif(nd, 0.2, 0.8))
}))
dat$y <- 1 + dat$x1 + rnorm(m)[dat$area] +
  rnorm(nrow(dat), sd = 0.5) + rnorm(nrow(dat), sd = sqrt(dat$vardir))

fit <- hb_twofold(y ~ x1, vardir = "vardir", domain = "area",
                  subarea = "subarea", data = dat)
summary(fit)
\end{lstlisting}

\subsection{Referensi kunci}\label{sec:hbtwofold-ref}

Blok roxygen memuat \citet{torabi2014small} (model dua-lipat),
\citet{rao2015small}, dan \citet{riebler2016intuitive}
(\code{R/hb\_twofold.R:105-112}); kerangka INLA mengacu ke
\citet{rue2009approximate} dan \citet{lindgren2015bayesian}.
